\section{相关工作}
\label{sec:7}

动态可组合性与若干已确立的研究领域存在交集。我们调研了最相关的工作路线，并将我们的贡献与每一条路线加以区分。

\subsection{效应与协效应系统}
\label{sec:7.1}

第 \S\ref{sec:2} 将效应与协效应作为支撑我们工作的理论支柱进行了综述。我们首先定位目前在工业实践中常见的单子效应系统，然后调研三个与 Cordis 相关的扩展效应与协效应方向的研究路线：将代数效应重新阐释为能力、赋予效应可逆语义，以及在单一分级规约下统一效应与协效应。

单子效应系统。一类库在现有通用语言的类型系统中编码效应，将其表示为运行时执行的单子值。Scala 中的 ZIO\cite{ref82} 将计算建模为 ZIO[R,E,A]，TypeScript 中的 Effect-TS\cite{ref83} 将其建模为 Effect<A,E,R>，这是一种泛型类型，其参数描述了计算的结果、类型化错误以及其上下文必须提供的服务；fp-ts 库\cite{ref84} 通过基于 Reader 的单子变换器编码了相同的错误和需求通道。两个特征将这些系统与 Cordis 区分开来。第一，追踪是以单子嵌入为代价获取的：程序只有在效应类型内部编写才能获得追踪，而 Cordis 将效应追踪作为普通宿主代码之上的叠加层。第二，需求通过解释来满足，即安装一个提供其操作的服务，而当该服务被撤回时，其操作所执行的内容仍然保留；Cordis 则将每个效应与一个逆操作配对，并随着提供者的到来和离去重新解析需求（\S\ref{subsec:3.1}、\S\ref{subsec:3.2}）。

代数效应作为能力。代数效应（\S\ref{sec:2.1}）使效应操作对类型系统可见。与我们工作最接近的扩展是 Brachthäuser 等人的 Effekt 语言，它将效应类型重新阐释为能力\cite{ref85,ref86}：效应类型表达的是计算从其上下文中需要什么，而非它可能产生哪些副作用。这一视角与我们的观点一样，将上下文视为能力的中介。Cordis 与 Effekt 在两个方面有所不同。(1) 在目的上，代数效应使效应可见以实现模块化解释，为一个操作赋予多种处理器语义，而 Cordis 使效应可见以实现追踪和回退，将每个上下文变换与一个逆操作配对。(2) 在设定上，Effekt 在类型层面静态规约效应，默认采用基于作用域的推理，其中能力是二等公民并被限制在其词法作用域内，通过装箱恢复一等使用，装箱通过在类型中追踪捕获的能力来解除该限制；Cordis 则在运行时规约效应，目标是组件移除时的完全资源恢复；\S\ref{sec:6.7} 讨论了在这个意义上将上下文视为二等公民的语言能提供什么。

可逆效应语义。一条平行的路线赋予效应可逆语义而非解释语义。Heunen 等人\cite{ref87}通过将 Hughes 的箭头适配为匕首箭头和逆箭头，在可逆设定中建模副作用，捕获了序列化和可变存储等操作可逆的效应。这是与我们的可逆效应最接近的形式化阐述：两者都将每个效应与撤销它的手段配对，而非通过处理器来满足。两者的区别在于可逆性所在的位置，以及它们对可逆性的需求程度。Heunen 等人在指称的、范畴论的设定中工作，其中可逆性是一种全局性质，由构造保证，因为每个计算都是可逆的，且逆是双向的并从范畴结构中恢复。Cordis 在运行时追踪逆操作，且要求更少：不要求整个计算可逆，而是每个原子效应允许一个单向逆操作，由调用者在应用点提供而非推导得出，由此任何复合效应的逆通过组合得出（\S\ref{subsec:3.1}）。

分级类型作为统一的效应与协效应。Orchard 等人\cite{ref88}提出了分级模态类型作为涵盖效应推理（通过分级单子）和协效应推理（通过分级余单子）的伞形概念，在 Granule 语言中实现，证明了单一类型系统可以同时追踪计算做什么和计算需要什么；更新的工作将协效应扩展到命令式的类 Java 语言\cite{ref89,ref90}和按值推入调用\cite{ref91}。所有这些都在类型层面操作：效应与协效应是在编译时对词法固定作用域检查的静态标注。我们的贡献与这一分析是正交的：我们将相同的两个概念提升为运行时机制，使得 Cordis 能够处理动态组合。时间回退和空间依赖随着已加载组件集合的演进而重新解析，而非在固定的程序文本上一次性确定。

\subsection{编程范式}
\label{sec:7.2}

\S\ref{subsubsec:3.3.3} 确立了上下文范式作为一种通过显式上下文中介效应与协效应的规约。两种已确立的范式值得明确比较：一种与我们有相同的术语，另一种与我们对横切关注点的处理方式相同。

面向上下文编程。COP\cite{ref92,ref93}为语言配备层——根据执行上下文在运行时激活和去激活的部分方法和类定义，使得行为无需基础代码命名其上下文依赖即可适应\cite{ref94}。COP 和 Cordis 在将上下文视为一等公民、运行时可变的实体以及在运行时动态激活和去激活行为方面一致，但这种相似只是名义上的。在 COP 中，"上下文"指的是周围的执行情境（例如位置、用户、模式），激活改变了动态作用域范围内的方法分派；层既不追踪它引发的副作用，也不回退它们，激活也不受依赖满足的约束。在 Cordis 中，上下文是中介效应与协效应的 $\Gamma_\infty$ 实体：激活运行组件的可逆效应并由响应式协效应满足驱动（\S\ref{subsec:3.2}），去激活则将其完全回退。COP 改变的是运行什么行为；Cordis 组合和回退的是组件安装了什么效应与依赖。它们的区别在于权衡。COP 将激活折叠到宿主语言的方法分派中，以语言特异性为代价获得动态作用域的层范围，而 Cordis 作为语言无关的叠加层，通过共享上下文响应式地解析激活。因此，Cordis 只能将 COP 全局的、值驱动的片段表达为协效应：在实现之间依赖于上下文的选择，但不能表达动态作用域的激活。

面向方面编程。AOP\cite{ref95,ref96}将横切关注点模块化为方面：一个切入点，它在基础程序中选定的连接点上进行量化，以及编入每个连接点的通知。Cordis 解决的是相同的上下文行为问题——这些行为否则会散布在各组件中——但其方面的对应物是协效应：一个共享的中介点，许多组件声明对其的依赖，使得横切行为可以在那里重塑而无需编辑它们中的任何一个。然后，两种范式在两个轴上有所不同。(1) 声明与不知情：AOP 切入点是不知情的和量化的，匹配其代码不知自己被通知的任意连接点，而 Cordis 将横切限制在每个组件声明的协效应，因此其范围恰好是所声明的表面。这带来了确定性和可追溯性：应用编排者可以在配置层检查和治理横切组件的内容，而无需阅读或分析其源代码，而 AOP 关注点只有通过对其进行量化的方面才能辨识。(2) 生命周期集成：Cordis 中的横切变更由组件的效应承载，在组件卸载时回退并响应式地传播到其依赖者，因此它是动态组合模型中的一次移动；动态 AOP 系统\cite{ref97,ref98}也可以在运行时织入和反织入，但作为独立操作，既不绑定到组件的生命周期，也不在被通知的代码中触发重新解析。

\subsection{时间可组合性}
\label{sec:7.3}

时间可组合性关注在运行程序中替换或移除组件时恢复其安装的效应。先前的方法按其处理离去组件的状态和效应的方式划分：将状态前向迁移到后继版本、通过开发者编写的清理恢复效应、在预先固定的作用域内自动逆转效应，或从运行时通过接口介入积累的记录中回收资源。

有状态前向迁移。一类广泛的系统通过在版本间前向迁移状态来替换运行程序中的组件而不停机。所有这些系统都遵循相同的时间规约：组件只有到达安全的、无交互的点时才能被替换。Kramer 和 Magee 将此准则确立为静止性\cite{ref51}，Vandewoude 等人后来将其放宽为干扰更小的宁静性\cite{ref52}；我们的滚动更新模式（\S\ref{sec:6.2}）通过在卸载提供者之前排空进行中的请求来强制执行该准则。然后，动态软件更新（DSU）通过手工编写的转换函数前向迁移状态：Hicks 等人面向 C 的通用 DSU\cite{ref99}、Stoyle 等人通过无关性分析实现的类型安全更新点\cite{ref100}，以及 Hayden 等人的 Kitsune\cite{ref101} 都将旧版本数据映射到新版本表示，继承堆对象、打开的文件和连接，同时重新初始化未迁移的部分。相同的规约扩展到持久状态：Overeem 等人\cite{ref102}通过手工编写的升级操作在模式版本之间转换运行中事件存储的数据，同时保持系统可用。Erlang/OTP\cite{ref15}在进程层面采取相同立场，通过 code\_change/3 迁移状态，并通过重启受监督的进程而非回退其效应来从故障中恢复；JavaScript 的热模块替换（例如 webpack\cite{ref46}、Vite\cite{ref47}）在模块层面做同样的事情，通过 module.hot 或 import.meta.hot API 在重新加载时前向传递状态。与 Cordis 的模块替换（\S\ref{sec:5.2}）相比，这些方法更优雅地迁移内存中的状态：Cordis 回退旧组件的已追踪效应并从干净状态重新应用新组件的效应，因此组件自身的内存中的状态不会在重新加载后保留，除非放置在更长寿命的依赖中，在可逆效应之上叠加 DSU 风格的前向迁移是未来的工作。尽管如此，Cordis 的方法在两个方面更为通用：它不需要 DSU 和 HMR 所需的手工编写迁移函数，而且它支持完全卸载组件并回收其资源，而不仅仅是就地更新。

开发者编写的恢复。第二类通过开发者手工编写的清理或补偿逻辑来恢复组件的效应。插件生命周期约定（例如 OSGi\cite{ref50}、Eclipse 的扩展点、IntelliJ 和 VSCode）将清理委托给开发者编写的卸载回调；命令模式\cite{ref103}将操作连同撤销方法封装以用于撤销/重做栈；Saga 模型\cite{ref49}将长生命周期事务结构化为每个步骤配有一个补偿动作的步骤序列；代数效应处理器可以附加在拆卸时运行终结器\cite{ref104}；事件溯源\cite{ref105}通过追加补偿事件而非执行逆操作来回退状态。在所有这些中，逆操作是一项未强制执行的职责，与操作分离，因此遗漏的逆操作会无声地泄漏资源（如 \S\ref{sec:1.2.1} 中经验性地记录的那样）。React 的 useEffect 钩子\cite{ref106}最接近于在结构上将效应与其逆操作配对，返回一个运行时在每次重新执行和卸载时调用的清理函数。其不足之处在于可组合性：钩子只能在组件或另一个钩子的顶层调用，不能在条件、循环或嵌套函数内部调用，且其效应体既不接受异步函数也不接受迭代器。因此，效应不能从其他效应组装或与控制流交错，没有可以从中推导复合逆操作的基础。Cordis 的效应没有这样的限制：它们是可以自由组合且可以异步运行的普通操作，只要求为每个原子效应手工编写逆操作，由此任何复合效应的逆通过组合得出，因此组装已有效应不需要编写任何逆操作。这种将每个效应与其逆操作的结构化配对使得完全恢复成为系统的不变量，而非开发者纪律的问题。

静态作用域的逆转。第三类通过构造自动逆转效应，但将逆转限制在预先固定的作用域内。软件事务性内存\cite{ref107,ref108}，源自硬件事务性内存\cite{ref109}，记录读/写日志，使得一组内存操作要么提交要么中止，将内存回滚到其事务前状态。可逆计算，从 Landauer 和 Bennett 的热力学分析\cite{ref110,ref111}到 Janus\cite{ref112}等可逆语言，更进一步使整个计算的每一步全局可逆。可逆进程演算将回溯构建到语义本身中：RCCS\cite{ref113}在每个进程旁携带一个记忆，并允许在其导致的历史因果等价时撤回一步，Phillips 和 Ulidowski\cite{ref114}统一地为 CCS、ACP 和 CSP 推导可逆算子，同时保持其前向操作语义。它们的因果一致性准则是 Cordis 恢复所遵循顺序的并发对应物——一个累加器以后进先出顺序应用组件自身的逆操作，\S\ref{sec:4.3.1} 的守卫延迟提供者的撤回直到其消费者已去激活（定理 63）。然而，其范围由语义固定，每个执行的动作都保持可撤销，而 Cordis 组件为每个原子效应提供一个逆操作，其累加器将上下文带回其组合开始的位置。线性类型\cite{ref115}、RAII\cite{ref4}和 Rust 的所有权系统\cite{ref61}将资源的释放绑定到词法区域。每种方法都静态固定了逆转的作用域和范围；相比之下，Cordis 不预先固定任何此类作用域：它在组件生命周期上回退任意的上下文操作，并将词法资源管理视为互补的，适合于单个组件内的局部资源。

介入式回收。第四类在组件本身不提供逆操作的情况下，通过在运行时控制的接口记录组件的获取来回收组件所获取的内容。Nooks\cite{ref77}封装了跨越 Linux 内核与其可加载扩展之间边界的每个调用，使得扩展接触的内核对象通过对象追踪器，其记录告诉恢复管理器在扩展故障时释放什么；影子驱动\cite{ref78}从另一侧接入相同的调用，记录请求和确定驱动程序状态的配置，以便重启的实例可以恢复到该状态。Akeso\cite{ref116}通过编译器插装获取记录，将内核执行划分为可嵌套的恢复域，记录其状态变更和跨线程依赖，并将故障请求与依赖于它的每个域一起回滚。因此，回收源于运行时维护的记录而非开发者记得编写的清理，这使得这一族系成为可逆效应最接近的系统级先例。它在与 Cordis 的词汇和范围上有所不同。平台固定了什么可以被记录，无论是每种内核对象类型的释放代码、每种驱动程序类的一个影子，还是每种插装分配器的一个逆操作，因此组件只能持有平台已知道如何释放的资源；相比之下，Cordis 组件引入自己的效应并为每个原子效应提供逆操作（\S\ref{subsec:3.1}）。回收同样受限于提交的请求或同一扩展的重启，而 Cordis 在组件的整个生命周期上回退并将移除传播到其依赖者，依赖者进而释放自己的效应（\S\ref{subsec:3.2}）。

\subsection{空间可组合性}
\label{sec:7.4}

空间可组合性关注组件对其他组件的依赖如何声明和绑定。先前的机制按绑定如何响应变化来划分：在初始化时一次性连接依赖、对整个组件的可用性做出响应，或在单个值的粒度上传播变化。

初始化时依赖连接。两种已确立的机制在初始化时将组件连接在一起。依赖注入框架\cite{ref38}（例如 Spring\cite{ref117}、Guice、Angular、Inversify）在初始化时将依赖注入组件，UI 框架上下文（例如 Vue.js 的 provide/inject 和 React 的 Context API）沿组件树传递依赖。一些支持动态作用域（例如 Spring 的 prototype/request 作用域、Angular 的分层注入器），但两者都不响应式地重新解析：当提供者在运行时被替换或移除时，现有依赖者既不去激活也不重新初始化，两者都不提供我们的组件状态机所提供的那种生命周期管理。Cordis 的响应式协效应（\S\ref{subsec:3.2}）提供了这一点：通知机制在满足谓词变化时触发生命周期转换。

可用性响应式组件模型。与我们的响应式协效应最接近的先例是对服务可用性做出响应。OSGi 的声明式服务和 iPOJO\cite{ref118,ref119}允许组件声明提供和需要的服务，运行时随着服务的出现和消失自动激活和去激活组件；iPOFO 的 Gravity 项目\cite{ref119}明确面向对变化的服务可用性的自主运行时适配，其 provide/require 模型直接预示了 Cordis 的 ctx.provide/ctx.get 模式。R-OSGi\cite{ref53}通过 RPC 将相同的抽象透明地扩展到分布式环境，将网络故障映射为服务撤回事件，\S\ref{sec:6.2} 将此模式作为 Cordis 模型的扩展进行讨论。所有这些系统通过去激活回调进行恢复，这在两个方面是有限的。第一，回调是手工编写的，因此资源安全依赖于开发者纪律，遗漏的回调会无声地泄漏。第二，回调是同步的：如果拆卸需要与离去的依赖进行异步交换，框架不提供等待它的协议，迫使对可能已经过时的引用进行阻塞等待。Cordis 的响应式协效应弥补了这两个差距：去激活回退依赖者的累积效应，其惯性 $\mathsf{Unloading}$ 状态（\S\ref{sec:4.3.3}）在处理进一步变更之前将异步拆卸运行至完成。

值级响应性。函数式响应式编程（FRP）\cite{ref120}及其现代形态（例如 SolidJS 中的信号\cite{ref121,ref122}、Vue 的响应性系统、Angular Signals）在值级粒度上传播变化：当信号变化时，派生计算被同步地或在调度器下重新求值\cite{ref123}。Cordis 的响应式协效应在组件级粒度上行动，添加了值级传播所不建模的异步生命周期语义。相同的粒度差异在一致性方面反过来同样成立：在一个轮次中，以依赖图固定的顺序传播，使得 FRP 能够要求没有派生计算读取更新和陈旧输入的混合，即无毛刺性\cite{ref124}，而 Cordis 没有轮次的对应物，编排动作逐个到达，只保证没有单个转换跨越其协效应的两次解析（定理 64）。两者是互补而非竞争的：Cordis 协效应本身可以携带响应式值，组件只在其实际消费的部分上更新，将组件级响应性细化为跨越两个层级的更细粒度响应式协效应。
